\documentclass[11pt]{article}
\usepackage[margin=1in]{geometry}
\usepackage{amsmath}
\usepackage{booktabs}
\usepackage{hyperref}

\title{Fixing Missing Initial Velocity for Newly Entering Inlet Droplets}
\author{Droplet rollout modeling notes}
\date{\today}

\begin{document}
\maketitle

\section{Problem}

The canonical droplet dataset stores each droplet state as
\[
    (x, y, v_x, v_y, c),
\]
where \(c\) is circularity. Velocities are currently computed by finite
differencing each track:
\[
    v_x(t) = x(t) - x(t-1), \qquad
    v_y(t) = y(t) - y(t-1).
\]
At the first visible frame of a track, there is no previous observation.
Therefore \(v_x\) and \(v_y\) are undefined and become NaN.

During model-window construction, non-finite values in the input history are
converted to zero. For newly entering inlet droplets this silently converts an
unknown inlet velocity into a stationary state. This is physically wrong: a
droplet appearing at the inlet is not stationary; it has just entered the field
of view.

\section{Observed Impact}

The failure case for track 3032 illustrates the problem. At the Markovian input
frame, the droplet has just appeared near the upper inlet, but its velocity is
missing:
\[
    y \approx 3.9,\qquad v_x = \mathrm{NaN},\qquad v_y = \mathrm{NaN}.
\]
After NaNs are converted to zero, the model receives an input resembling a
stationary droplet near the inlet. The true future motion, however, rapidly
reaches approximately \(11\) pixels per frame in the \(y\) direction.

In validation diagnostics, this start-state issue affected 110 evaluable
trajectories, or about \(0.62\%\) of all evaluable validation trajectories. It
was highly concentrated at the inlet: 109 of 110 affected cases had input
\(y < 10\) pixels. Although rare by sample count, these cases were strongly
associated with severe channel-admissibility failures in the geometry-aware
Markovian rollout.

\section{Scientific Interpretation}

This is a data representation issue, not a model architecture issue. The
observed state is intended to include velocity. For newly entering droplets,
the velocity is not measured at the first visible frame, but the current
preprocessing turns that missing value into a literal zero. The model is
therefore asked to interpret an ambiguous and physically misleading state.

The issue is especially important for Markovian models, because they receive
only the current state. A history-based model can often infer motion from
previous frames when they exist, but a one-frame Markovian model cannot infer
an inlet velocity from a single position if the velocity channel has been
zeroed.

\section{Required Correction}

The correction should be applied once, when constructing the canonical dataset,
before train/validation/test window construction and before any model-specific
data loading. This ensures that all models consume the same corrected state
representation.

Let \(\mathcal{I}\) denote finite velocity observations in the incoming inlet
region. The inlet prior velocity is
\[
    \bar{v}_x^{\mathrm{inlet}}
    =
    \frac{1}{|\mathcal{I}|}
    \sum_{i \in \mathcal{I}} v_{x,i},
    \qquad
    \bar{v}_y^{\mathrm{inlet}}
    =
    \frac{1}{|\mathcal{I}|}
    \sum_{i \in \mathcal{I}} v_{y,i}.
\]

For newly visible inlet droplets only, replace missing velocity components:
\[
    v_x \leftarrow \bar{v}_x^{\mathrm{inlet}}
    \quad \mathrm{if}\quad v_x=\mathrm{NaN},
\]
\[
    v_y \leftarrow \bar{v}_y^{\mathrm{inlet}}
    \quad \mathrm{if}\quad v_y=\mathrm{NaN}.
\]

The patch must not be applied to arbitrary missing velocities elsewhere in the
channel. The intended criterion is:
\begin{enumerate}
    \item the row is the first visible frame of its track,
    \item the droplet is in the incoming inlet region,
    \item at least one of \(v_x, v_y\) is non-finite.
\end{enumerate}

\section{Diagnostics}

The canonical dataset build should report:
\begin{itemize}
    \item the inlet-region criterion,
    \item \(\bar{v}_x^{\mathrm{inlet}}\),
    \item \(\bar{v}_y^{\mathrm{inlet}}\),
    \item number of droplet rows patched,
    \item number of velocity values patched,
    \item number of affected unique tracks.
\end{itemize}

\section{Expected Result}

After rebuilding the canonical dataset, first-visible inlet droplets should no
longer enter model inputs with zeroed velocities. History-20, Markovian, and
geometry-aware Markovian models should all train and evaluate from the same
corrected canonical state tensor. No model architecture, rollout logic,
boundary injection, loss function, geometry mask, or ellipse construction
should change.

\end{document}
